\subsection{Magnus 群与自由群}

定理 1. 设 \(r\) 是 \(\mathbf{X}\) 到 \(\hat{\mathbf{A}}(\mathbf{X})\) 中的一个映射，使得 \(\omega(r(x)) \geqslant 2\) （对一切 \(x \in \mathbf{X}\) ）。唯一的同态 \(g\) ——它把自由群 \(\mathrm{F}(\mathbf{X})\) 映到 Magnus 群 \(\Gamma(\mathbf{X})\) 中，并且使得 \(g(x) = 1 + x + r(x)\) （对一切 \(x \in \mathbf{X}\) ）——是单射。

我们先证明三个引理。

引理 2. 设 \(n\) 是一个非零有理整数。在形式幂级数环 \(\mathbf{K}[[t]]\) 中我们记 \((1 + t)^n = \sum_{j \geqslant 0} c_{j,n} t^j\) 。存在一个整数 \(j \geqslant 1\) 使得 \(c_{j,n} \neq 0\) 。

若 \(n > 0\) ，则由二项式公式有 \(c_{n,n} = 1\) 。

设 \(n < 0\) ，并令 \(m = -n\) 。若 \(c_{j,n} = 0\) （对一切 \(j \geqslant 1\) ），则 \((1 + t)^n = 1\) ，从而取逆便得 \((1 + t)^m = 1\) ，这与公式 \(c_{m,m} = 1\) 矛盾。

引理 3. 设 \(x_1, \ldots, x_s\) 是 \(\mathbf{X}\) 的元素，满足 \(s \geqslant 1\) ，并且 \(x_i \neq x_{i+1}\) （对 \(1 \leqslant i \leqslant s-1\) ）；设 \(n_1, \ldots, n_s\) 是非零有理整数。那么元素 \(\prod_{i=1}^{s} (1 + x_i)^{n_i}\) （属于 \(\hat{\mathbf{A}}(\mathbf{X})\) ）是 \(\neq 1\) 。

设 \(m\) 是 \(K\) 的一个极大理想， \(k\) 是域 \(K/m\) ；设 \(p: \hat{A}_K(X) \to \hat{A}_k(X)\) 是唯一的连续幺 \(K\) -代数同态，使得 \(p(x) = x\) （对 \(x \in X\) ）(no. 1， 命题 1)。只需证明 \(p\left(\prod (1 + x_i)^{n_i}\right) \neq 1\) ，问题就化为 \(K\) 是域的情形。

用引理 2 的记号：

\[
\prod_ {i = 1} ^ {s} (1 + x _ {i}) ^ {n _ {i}} = \sum_ {b _ {i} \geqslant 0} c _ {b _ {1}, n _ {1}} \dots c _ {b _ {s}, n _ {s}} x _ {1} ^ {b _ {1}} \dots x _ {s} ^ {b _ {s}}.
\]

由引理 2，存在整数 \(a_i > 0\) 使得 \(c_{a_i, n_i} \neq 0 (1 \leqslant i \leqslant s)\) 。由《代数》第 I 章 § 7 no. 4 命题 6，没有单项式 \(x_1^{b_1} \ldots x_s^{b_s}\) （满足 \(b_{i} \geqslant 0\) 且 \((b_{1}, \ldots, b_{s}) \neq (a_{1}, \ldots, a_{s})\) ）能够等于 \(x_{1}^{a_{1}} \ldots x_{s}^{a_{s}}\) 。由此得出 \(x_{1}^{a_{1}} \ldots x_{s}^{a_{s}}\) 在 \(\prod_{i=1}^{s} (1 + x_{i})^{n_{i}}\) 中的系数是 \(c_{a_{1}, n_{1}} \ldots c_{a_{s}, n_{s}} \neq 0\) ，这就蕴涵所述结果。

引理 4. 设 \(\sigma\) 是 \(\hat{\mathbf{A}}(\mathbf{X})\) 的连续自同态，使得 \(\sigma(x) = x + r(x)\) （对 \(x \in \mathbf{X}\) ）(no. 1， 命题 1)。那么 \(\sigma\) 是一个自同构，并且 \(\sigma(\hat{\mathbf{A}}_m(\mathbf{X})) = \hat{\mathbf{A}}_m(\mathbf{X})\) （对一切 \(m \in \mathbf{N}\) ）。

\(\sigma(x) \equiv x \bmod. \hat{A}_2(X)\) （对 \(x \in X\) ），从而对 \(n \geqslant 1\) 与 \(x_1, \ldots, x_n\) （属于 \(X\) ），

\[
\sigma (x _ {1}) \dots \sigma (x _ {n}) \equiv x _ {1} \dots x _ {n} \mod \hat {\mathrm{A}} _ {n + 1} (\mathrm{X});
\]

由线性性推出 \(\sigma(a) \equiv a\) 模 \(\hat{\mathbf{A}}_{n+1}(\mathbf{X})\) 成立（对一切 \(a \in \mathbf{A}^n(\mathbf{X})\) ），特别地 \(\sigma(\mathbf{A}^n(\mathbf{X})) \subset \hat{\mathbf{A}}_n(\mathbf{X})\) 。由此得出 \(\sigma(\hat{\mathbf{A}}^m(\mathbf{X})) \subset \hat{\mathbf{A}}_n(\mathbf{X})\) （对 \(m \geqslant n\) ），从而 \(\sigma(\hat{\mathbf{A}}_n(\mathbf{X})) \subset \hat{\mathbf{A}}_n(\mathbf{X})\) 。换句话说， \(\sigma\) 与滤过 \((\hat{\mathbf{A}}_m(\mathbf{X}))\) （\(\mathbf{A}(\mathbf{X})\) 上的）相容，并且它在相伴分次环上的限制是恒同映射。因此 \(\sigma\) 是双射 (《交换代数》， 第 III 章， § 2， no. 8， 定理 1 的推论 3)。

最后我们证明定理 1。设 \(w \neq 1\) 是 \(\mathrm{F}(\mathbf{X})\) 的一个元素。由《代数》第 I 章 § 7 no. 5 命题 7，存在 \(x_{1}, \ldots, x_{s}\) （属于 \(\mathbf{X}\) ）以及非零有理整数 \(n_{1}, \ldots, n_{s}\) ，使得 \(s \geqslant 1\) 、 \(x_{i} \neq x_{i+1}\) ( \(1 \leqslant i \leqslant s - 1\) ) 并且

\[
w = x _ {1} ^ {n _ {1}} \dots x _ {s} ^ {n _ {s}}.
\]

用引理 4 的记号，

\[
g (w) = \prod (1 + \sigma (x _ {i})) ^ {n _ {i}} = \sigma (\prod (1 + x _ {i}) ^ {n _ {i}}),
\]

因而由引理 3 与引理 4 得 \(g(w) \neq 1\) 。

